\section{Review Scope and Analytical Lens}\label{sec:scope}

This narrative review draws on 141 references: 130 scholarly papers and book chapters, one book, and 10 standards, technical reports, or other official technical references. The literature is in English and mainly covers publications from 2016--2026, supplemented by earlier foundational work on reputation, distributed-system limits, provenance, and consensus. The scholarly literature includes original studies, reviews, and papers describing datasets, benchmarks, and simulation platforms.

The references cover six related themes: UAV blockchain and swarm security; maritime and underwater communication and security; heterogeneous UAV--USV--UUV cooperation; IoT trust and reputation; provenance, attestation, and zero trust; and permissioned-ledger architecture and governance. Unmanned-system studies establish the operational constraints of air, surface, and underwater missions. General IoT, cybersecurity, and distributed-systems research provides mechanisms for authentication, evidence appraisal, data sharing, and coordination. Standards and official specifications supply architectural definitions and interoperability requirements, while datasets and simulation platforms provide resources for evaluation.

The analysis compares mechanisms by the object being assessed, the assurance claim they support, and the mission stage at which they operate. These three dimensions form the Object--Assurance--Lifecycle (OAL) framework. Each comparison considers the supporting evidence and the network, resource, threat, and governance assumptions that condition its use. Transfer to a maritime mission depends on whether those assumptions hold in the target environment. For example, a consensus mechanism evaluated on shore servers requires separate assessment under the energy limits and intermittent acoustic connectivity of a UUV. This perspective connects the literature to the dependency analysis, CTG reasoning loop, and evaluation agenda developed in later sections.

Table~\ref{tab:scope} summarizes the review scope. UAVs, USVs, and underwater vehicles form the operational focus, with satellite, shore, cloud, and edge systems providing supporting infrastructure. The lifecycle extends from admission and sensing through mission decisions to audit, revocation, and recovery, encompassing both operational trust and accountability across participating organizations.

\begin{table}[!htbp]
\caption{Review scope and analytical focus.\label{tab:scope}}
\begin{tabularx}{\textwidth}{@{}P{29mm}YY@{}}
\toprule
\textbf{Dimension} & \textbf{Analytical focus} & \textbf{Complementary context}\\
\midrule
Physical domains & UAV, USV, UUV/AUV/ROV; buoys and cross-media gateways & Ground-vehicle analogies and platform structural-certification requirements\\
Infrastructure & Satellite, high-altitude, shore, cloud, edge, and consortium services & Partition-aware operation and explicit infrastructure trust assumptions\\
Mission settings & Observation, inspection, search and rescue, environmental monitoring, logistics, and publicly documented defense-support contexts & Publicly reportable assurance semantics, failure modes, and evaluation methods\\
Trust mechanisms & Identity, authorization, attestation, provenance, dynamic inference, policy, ledger, audit, and governance & Resource requirements, evidence freshness, uncertainty, and conditions for cross-domain use\\
Evidence & Original studies, reviews, a scholarly book, standards, and official technical references & Datasets, benchmarks, simulation platforms, and foundational work from related fields\\
Review output & Conceptual framework, taxonomy, design boundaries, evidence profile, and research propositions & Empirical comparison, formal verification, and deployment certification as evaluation stages\\
\bottomrule
\end{tabularx}
\end{table}
